\documentclass[../candidatura.tex]{subfiles}
\graphicspath{{\subfix{../imagenes/}}}
\begin{document}

\section{Síntesis sonora}

\avisoesbozo

Síntesis en español viene del latín synthĕsis \parencite{wiktionary_synthesis_sinfecha}, y este a su vez del griego \foreignlanguage{greek}{σύνθεσις}. Esta palabra tiene dos raíces ``sin'' y ``tesis''.

La primera se traduce como en conjunto, la segunda como armar, rejuntar, con lo que síntesis evoca la composición de un todo a partir de otras partes.

Si constelamos o caleidoscopamos estas partículas semánticas, y contrastamos con otras acepciones de síntesis en español, también emergen hálitos de resumen, de artificial, de suma, de integración, de proceso químico.

Desde la adopción de circuitos eléctricos para la creación de instrumentos sonoros en el siglo XX, se adoptó el nombre de síntesis para referirse a la creación de sonidos a partir de otras fuentes no acústicas.

Existen subvariantes como síntesis aditiva, substractiva, granular, por modulación de frecuencia, por modulación de anillo, entre otras.

Emergieron en el siglo XX dos grandes escuelas de pensamiento sobre nuevas formas de sonar, en Francia la Musique concrète, en Alemania Elektronische Musik.

Ambos paradigmas con sus acercamientos y filosofías de producción sonora continúan vigentes hoy en distintos formatos, desde inspirar el lenguaje usado para describir la síntesis, como también en las arquitecturas de los hardware y software para producir sonido.

En este capítulo se abordarán dos tradiciones europeas de síntesis sonora, incluyendo su influencia en instrumentos actualmente en producción.


\subsection{Musique concrète}

Muchas técnicas de la música concreta francesa han resonado con distintos módulos diseñados por Tom Erbe, y su empresa soundHack en colaboración con Make Noise Music. 

\subsection{Elektronische Musik}

La tradición de usar distintos módulos electrónicos originalmente diseñados para laboratorios de radio, entre otros, ha impactado en el desarrollo de paradigmas de síntesis modular con cableado.

Hay ecos de esta práctica en la programación por nodos en entornos como Max, TouchDesigner, y Unreal Engine, además del nexo con los sintetizadores modulares de empresas como ARP o los estándares Eurorack contemporáneos.

\newpage

\end{document}
